\documentclass[10pt]{article}
\providecommand{\FIGDIR}{fig}
\newcommand{\DIGESTDATE}{1 August 2026}
% ===================== PM Research Watch - shared preamble =====================
\usepackage[T1]{fontenc}
\usepackage[utf8]{inputenc}
\usepackage{amsmath,amssymb}
\usepackage{tgpagella}
\usepackage{tgheros}
\renewcommand{\sfdefault}{qhv}
\renewcommand{\ttdefault}{lmtt}
\usepackage{microtype}
\usepackage[english]{babel}

\usepackage[a4paper,top=17mm,bottom=17mm,left=17mm,right=17mm,
            headsep=6mm,footskip=10mm]{geometry}
\usepackage{graphicx}
\usepackage{xcolor}
\usepackage{colortbl}
\usepackage{tikz}
\usetikzlibrary{calc}
\usepackage[most]{tcolorbox}
\usepackage{titlesec}
\usepackage{fancyhdr}
\usepackage{booktabs}
\usepackage{array}
\usepackage{longtable}
\usepackage{enumitem}
\usepackage{multicol}
\usepackage{ragged2e}
\usepackage[hidelinks]{hyperref}

% ------------------------------------------------------------------ palette
\definecolor{Ink}    {HTML}{14262E}
\definecolor{Deep}   {HTML}{0F4C5C}
\definecolor{Teal}   {HTML}{1F7A8C}
\definecolor{Sky}    {HTML}{4FA3C4}
\definecolor{Sage}   {HTML}{6FA287}
\definecolor{Amber}  {HTML}{E3A008}
\definecolor{Clay}   {HTML}{B5643C}
\definecolor{Coral}  {HTML}{D1495B}
\definecolor{Violet} {HTML}{6A5B9E}
\definecolor{Slate}  {HTML}{4A6572}
\definecolor{Paper}  {HTML}{FBF9F5}
\definecolor{Mist}   {HTML}{EEEAE2}
\definecolor{Rule}   {HTML}{D9D3C9}

\pagecolor{Paper}
\color{Ink}

\hypersetup{colorlinks=true, linkcolor=Deep, urlcolor=Teal, citecolor=Teal,
            pdftitle={PM Research Watch}, pdfauthor={Automated daily digest}}

% ------------------------------------------------------------------ headings
\titleformat{\section}
  {\sffamily\bfseries\large\color{Deep}}{}{0pt}
  {}[{\vspace{1.2mm}\color{Amber}\titlerule[1.1pt]}]
\titlespacing*{\section}{0pt}{5.5mm}{2.8mm}

\titleformat{\subsection}
  {\sffamily\bfseries\normalsize\color{Teal}}{}{0pt}{}
\titlespacing*{\subsection}{0pt}{3.6mm}{1.4mm}

\setlist[itemize]{leftmargin=4.2mm, itemsep=1.1mm, topsep=1.3mm, parsep=0pt,
                  label=\textcolor{Teal}{\footnotesize$\blacktriangleright$}}

% ------------------------------------------------------------------ header/footer
\pagestyle{fancy}
\fancyhf{}
\renewcommand{\headrulewidth}{0pt}
\renewcommand{\footrulewidth}{0pt}
\fancyhead[L]{\sffamily\scriptsize\color{Slate}PM RESEARCH WATCH\;\textcolor{Rule}{\textbar}\;\DIGESTDATE}
\fancyhead[R]{\sffamily\scriptsize\color{Slate}Particulate matter monitoring \& health effects}
\fancyfoot[C]{\sffamily\scriptsize\color{Slate}\thepage}

% ------------------------------------------------------------------ components
% metric tile
\newcommand{\metric}[3]{%
  \begin{tcolorbox}[enhanced, colback=white, colframe=Rule, boxrule=0.5pt,
    arc=1.6mm, left=2.2mm, right=2.2mm, top=1.8mm, bottom=1.8mm,
    borderline west={1.7mm}{0pt}{#3}, nobeforeafter, valign=center,
    height=17mm, width=\linewidth]
    \RaggedRight
    {\sffamily\bfseries\fontsize{17}{17}\selectfont\color{#3}#1}\\[0.5mm]
    {\sffamily\fontsize{6.5}{7.4}\selectfont\color{Slate}#2}
  \end{tcolorbox}}

% pull-quote / key insight box
\newtcolorbox{keybox}[1][]{enhanced, colback=Mist, colframe=Mist,
  arc=1.6mm, left=3.2mm, right=3.2mm, top=2.4mm, bottom=2.4mm,
  borderline west={1.7mm}{0pt}{Amber}, #1}

% subtopic band
\newcommand{\band}[2]{%
  \vspace{3.2mm}\noindent
  \begin{tcolorbox}[enhanced, colback=#2, colframe=#2, arc=1.2mm,
    left=2.6mm, right=2.6mm, top=1.1mm, bottom=1.1mm, boxrule=0pt,
    width=\linewidth, nobeforeafter]
    {\sffamily\bfseries\small\color{white}#1}
  \end{tcolorbox}\vspace{1.6mm}}

% one paper entry: {short cite}{title}{journal}{doi}{body}
\newcommand{\paperentry}[5]{%
  \par\vspace{1.9mm}\noindent
  {\sffamily\bfseries\footnotesize\color{Ink}#1}\;%
  {\sffamily\fontsize{7}{8}\selectfont\color{Slate}#3}\\[0.4mm]
  {\footnotesize\itshape\color{Deep}#2}\\[0.6mm]
  {\footnotesize #4}\\[0.3mm]
  {\fontsize{6.6}{7.6}\selectfont\color{Slate}\texttt{doi:}\,\href{https://doi.org/#5}{\texttt{\detokenize{#5}}}}
  \par}

% caption style
\newcommand{\figcap}[1]{\par\vspace{1.0mm}
  {\fontsize{7.2}{8.6}\selectfont\color{Slate}\textsf{\textbf{Fig.}}\; #1\par}\vspace{1.5mm}}

\providecommand{\permil}{\textperthousand}

\setlength{\parindent}{0pt}
\setlength{\parskip}{0pt}
\renewcommand{\arraystretch}{1.18}


\begin{document}

% ==================================================================== MASTHEAD
\thispagestyle{empty}
\begin{tikzpicture}[remember picture, overlay]
  \fill[Deep] (current page.north west) rectangle
      ($(current page.north east)+(0,-67mm)$);
  \fill[Amber] ($(current page.north west)+(0,-67mm)$) rectangle
      ($(current page.north east)+(0,-68.4mm)$);
  \foreach \x/\y/\r in {22/12/0.9, 41/26/1.5, 63/9/0.7, 88/31/1.1, 112/17/1.9,
                        138/29/0.8, 159/13/1.3, 178/54/1.0, 196/21/1.6,
                        30/44/0.6, 74/57/1.2, 125/48/0.9, 168/6/0.8, 191/59/0.7,
                        52/18/1.0, 100/60/0.6, 148/20/1.1, 66/49/1.4, 205/45/0.8}
    \fill[Sky, opacity=0.30] ($(current page.north west)+(\x mm,-\y mm)$) circle (\r mm);
\end{tikzpicture}

\vspace*{4mm}
{\sffamily\fontsize{7.6}{9}\selectfont\color{Sky}%
 AUTOMATED DAILY LITERATURE DIGEST\;\textbullet\;PhD RESEARCH WATCH\par}
\vspace{2.2mm}
{\sffamily\bfseries\fontsize{27}{29}\selectfont\color{white}%
 Particulate Matter\par}
{\sffamily\fontsize{27}{29}\selectfont\color{Sky}%
 Monitoring \& Health Effects\par}
\vspace{5.0mm}
{\sffamily\fontsize{8.4}{10}\selectfont\color{white}%
 Issue 2026-08-01\;\;\textcolor{Amber}{\textbar}\;\;Entry window 1 August 2026,
 plus first sweep of the new by-journal leg\;\;%
 \textcolor{Amber}{\textbar}\;\;25 records in scope\par}
\vspace{18mm}

% -------------------------------------------------------------- metric strip
\noindent
\begin{minipage}[t]{0.190\linewidth}\metric{25}{RECORDS IN SCOPE\\(8 SUBTOPICS)}{Deep}\end{minipage}\hfill
\begin{minipage}[t]{0.190\linewidth}\metric{9}{SENSING /\\INSTRUMENTATION}{Teal}\end{minipage}\hfill
\begin{minipage}[t]{0.190\linewidth}\metric{17}{FULL ENTRIES\\(8 METADATA-ONLY)}{Sky}\end{minipage}\hfill
\begin{minipage}[t]{0.190\linewidth}\metric{5}{RECORDS FROM CHINA\\(20\% OF CORPUS)}{Amber}\end{minipage}\hfill
\begin{minipage}[t]{0.190\linewidth}\metric{2}{EFFECT ESTIMATES\\EXTRACTED}{Clay}\end{minipage}

\vspace{6mm}

% -------------------------------------------------------------- signal of day
\section*{\sffamily\bfseries\color{Deep}Signal of the day}
\vspace{-1.5mm}{\color{Amber}\rule{\linewidth}{1.1pt}}\vspace{2.5mm}

\begin{keybox}
{\sffamily\bfseries\small\color{Deep}Incense arrives from both ends on the same day: a
birth cohort finds a dose--response for small-for-gestational-age, and a chamber study
shows why a mass-based monitor would have missed the exposure entirely.}\\[1.4mm]
\footnotesize
Pai et al.\ (\textit{Pediatr Neonatol}) analysed \textbf{18{,}666 mother--infant pairs}
in the Taiwan Birth Cohort Study with household incense exposure categorised as none,
occasional or daily. Among term infants, SGA odds were \textbf{aOR 1.16 (95\% CI
1.00--1.33)} for occasional and \textbf{1.25 (1.10--1.41)} for daily burning --- an
ordered dose--response, robust to adjustment for other indoor sources, maternal medical
history, gestational weight gain and secondhand smoke, with no association for preterm
birth. Exposure is self-reported and unquantified, which is the study's real limit: it
establishes that the practice matters without saying what dose does.

Zhou et al.\ (\textit{Atmos.\ Chem.\ Phys.}) supply the missing half. Chamber
characterisation of Arabian incense (Bakhoor) under charcoal heating gives emission rates
of \textbf{670--1690\,$\mu$g\,min$^{-1}$\,g$^{-1}$} by mass and
\textbf{(6--7)$\times$10$^{11}$ particles\,min$^{-1}$\,g$^{-1}$} by number, with
\textbf{ultrafine particles contributing 70--75\% of total particle number} and UFP
emission rates exceeding sidestream cigarette smoke. Water-soluble DTT activity was
$\sim$32\,pmol\,min$^{-1}$\,$\mu$g$^{-1}$ --- modestly \emph{below} cigarette particles ---
yet the intracellular oxidative-stress response was \emph{stronger}, and ozone ageing
raised oxidative potential for both sources.

Read together the operational point is unambiguous. The health signal sits in a source
whose emission is dominated by particle \emph{number} in the ultrafine range, whose
toxicity ranking inverts between acellular and cellular assays, and whose burden is
therefore invisible to the PM$_{2.5}$-mass instrument that every relevant epidemiological
study --- including Pai's --- would have deployed. Number-resolved, low-cost indoor
sensing is not a refinement here; it is the only way to put an exposure metric on the
association.
\end{keybox}

\vspace{3mm}

% -------------------------------------------------------------- exec brief
\section*{\sffamily\bfseries\color{Deep}What changed today}
\vspace{-1.5mm}{\color{Amber}\rule{\linewidth}{1.1pt}}\vspace{2.5mm}

\footnotesize
\begin{itemize}
\item \textbf{The sensing blackout is over, and it was a harvest defect exactly as
diagnosed.} The 31 July issue recorded a second consecutive window with zero sensing
records and prescribed the fix: a Crossref by-journal feed keyed on \texttt{created}.
That leg was built and run today across eight journals PubMed does not index or indexes
late. It returned \textbf{55 deposits} for 26 July--1 August, of which 12 were
PM-relevant. Today's issue carries \textbf{9 sensing/instrumentation records} against
0, 0 and 1 in the three preceding issues. The defect was never a real absence in the
literature.

\item \textbf{Calibration papers now agree on what breaks, and it is not the sensor.}
Kumpika et al.\ reach $R^2 = 0.974$ against a BAM with a third-order humidity correction
\emph{and} report that duty-cycling a PMS7003 to 20\,min/day slows drift relative to
continuous operation; Qian et al.\ halve NRMSE with AutoML but only by splitting models
across clean-air and pollution-event concentration ranges; Le et al.\ trace a
ceilometer's signal loss to window fogging from insufficient firmware-controlled heating,
and show the internal calibration factor stops compensating once laser power falls below
40\%. Three different instrument classes, one conclusion: the dominant error terms are
environmental and ageing-related, and they are not on the manufacturer's specification
sheet.

\item \textbf{Exposure-model methodology got two hard results on the same failure mode.}
Saeedi et al.\ show land-use regression for mobile ultrafine particles collapsing from
\textbf{$\sim$217\% average percentage error under random cross-validation to
$\sim$79\% under spatiotemporal CV with forward feature selection} --- i.e.\ the standard
validation protocol was reporting a model that could not predict. Blanco et al.\ then
resample on-road monitoring designs against a 5{,}283-person cohort and find
health-effect estimates attenuate under every alternative design, with short 4-visit
campaigns the most unstable. Both papers say the same thing: reported LUR performance is
a property of the validation scheme, not of the exposure surface.

\item \textbf{Two independent suicide records surfaced in one window.} Daoud et al.\
report a fixed-effects panel association across UK local authorities
($\beta = 0.221$, $p = 0.007$) with a U-shaped form turning at
$\approx$7.8\,$\mu$g/m$^3$; Cai et al.\ (\textit{Atmos.\ Environ.}, metadata only) report
PM$_{2.5}$ as the dominant pollutant in a province-wide case-crossover among people
living with HIV/AIDS in China. Neither is strong evidence alone --- Daoud's is ecological
with no time-varying socioeconomic controls --- but the convergence marks a subtopic that
is accumulating.

\item \textbf{Source attribution keeps demoting the sector everyone regulates first.}
Zhou et al.'s 2022 India WRF-Chem run (annual bias $0.2 \pm 16.9$\,$\mu$g/m$^3$ across
288 sites, population-weighted mean 47.4\,$\mu$g/m$^3$) finds \textbf{transboundary
transport contributing 27\%} --- more than any single domestic sector --- with industry at
18\% now ahead of residential at 15\% nationally. The residential decline is a real
policy result; the transboundary share is a jurisdictional problem no national
instrument reaches.

\item \textbf{Oxidative potential became a modelled quantity.} Mishra et al.'s KM-OP
kinetic model reproduces measured OP in Grenoble, Paris and London at $R^2 > 0.75$ with
$<15\%$ bias in 4 of 6 datasets, identifies organic aerosol as the strongest contributor
to antioxidant-based assays, and separates Cu (dominates H$_2$O$_2$) from the
\textbullet OH production that actually tracks the assays. This is the first record in
this watch that lets OP be predicted from composition rather than measured per filter.
\end{itemize}

\clearpage

% ==================================================================== FIGURES
\section*{\sffamily\bfseries\color{Deep}Corpus analytics}
\vspace{-1.5mm}{\color{Amber}\rule{\linewidth}{1.1pt}}\vspace{2.5mm}

\noindent\includegraphics[width=\linewidth]{\FIGDIR/f1_subtopics.png}
\figcap{Subtopic assignment of the 25 in-scope records, single-label by dominant
contribution. \textit{Sensing, forecasting \& instrumentation} is the largest cluster at
9 records --- the first time it has led an issue --- and is entirely a product of the
by-journal leg activated today. Exposure assessment follows at 4. The health clusters are
correspondingly thin because the PubMed entry window covers a single day (1 August) while
the sensing leg back-swept a week.}

\vspace{2mm}
\noindent\includegraphics[width=\linewidth]{\FIGDIR/f2_design_endpoint.png}
\figcap{Left --- study architecture. Eight records are metadata-only (Elsevier titles
whose abstracts are deposited in neither Crossref nor Europe PMC at this point in the
publication cycle) and are given their own slice rather than an inferred design; the
remaining 17 span modelling/inventory, measurement campaign, review, cohort,
cross-sectional, ecological and experimental architectures.
Right --- health endpoint by organ system. \textbf{Ten records carry no human health
endpoint at all and eight more were not opened}, which is what an
instrumentation-weighted issue looks like. The endpoints that do appear are
neuro/cognitive (2), neuro/mental health (1), reproductive (2), infection risk (1) and
an in-vitro oxidative-stress proxy (1).}

\vspace{2mm}
\noindent\includegraphics[width=\linewidth]{\FIGDIR/f3_metric_geography.png}
\figcap{Left --- particle metric. \textbf{Ultrafine particle number appears in three
records today} (Bakhoor emissions, Toronto mobile LUR, Seattle on-road design) against a
recent baseline of zero or one --- the direct consequence of reading the aerosol-science
journals rather than only the biomedical ones. Right --- geographic provenance. China
leads at 5 (20\%), well below its usual share, and chamber and simulation studies carry
no national attribution.}

\vspace{2mm}
\noindent\includegraphics[width=\linewidth]{\FIGDIR/f6_lifecourse.png}
\figcap{Life-course windows addressed; counts are non-exclusive and only records with an
identifiable human population are counted. Working age leads at 3 (suicide mortality,
hospital bioaerosol exposure, school indoor exposure), with the in-utero window at 2.
Adolescence is empty again.}

\vspace{2mm}
\noindent\includegraphics[width=\linewidth]{\FIGDIR/f4_heatmap.png}
\figcap{Subtopic $\times$ study-architecture cross-tabulation. The structural feature to
note is that the sensing row is now populated across four distinct architectures ---
field co-location, chamber, network deployment and simulation --- rather than the single
forecasting-model cell it collapsed to in previous issues. The health rows remain
observational-only; no record today links an instrument characterisation to a health
endpoint within the same study.}

\vspace{2mm}
\noindent\includegraphics[width=\linewidth]{\FIGDIR/f5_forest.png}
\figcap{The two quantitative effect estimates in the corpus, both from Pai et al.'s
incense--SGA analysis, on a common logarithmic axis. \textbf{Two estimates is an
unusually thin forest and reflects the issue's composition rather than the literature}:
instrumentation and source-attribution papers report performance metrics, not risk
ratios, and the eight metadata-only records were not opened. Zhang et al.\ report odds
increases of 19--43\% per IQR across six PM$_{2.5}$ components but give no per-component
interval in the abstract; Bello et al.\ report $-$6.6\,g birth weight per
1\,$\mu$g/m$^3$ in the most vulnerable stratum with no published CI; Daoud et al.\ report
a panel coefficient. None of the three is plottable without inventing an interval, so
none was plotted.}

\clearpage

% ==================================================================== DIGEST
\section*{\sffamily\bfseries\color{Deep}Paper-level digest}
\vspace{-1.5mm}{\color{Amber}\rule{\linewidth}{1.1pt}}\vspace{1mm}

{\fontsize{7.4}{9}\selectfont\color{Slate}Ordered by cluster. Each entry gives the
methodological core and the finding that matters, not an abstract paraphrase. Eight
records surfaced without a retrievable abstract are listed separately at the end on
metadata alone --- they are not summarised, because there is nothing to summarise from.\par}

\band{Sensing, forecasting \& instrumentation\hfill 5 full entries}{Teal}

\paperentry{Kumpika et al.\ 2026}
{Correction of humidity and sensor ageing effects in low-cost PM$_{2.5}$ light-scattering
sensors}
{Environ Pollut Bioavailab}
{Field campaign in Chiang Mai with PMS7003 units across a wide relative-humidity range,
validated against a government BAM and cross-checked with dry-air measurement through a
diffusion dryer. Polynomial corrections up to third order were fitted;
\textbf{agreement with BAM reached $R^2 = 0.9737$} after correction, with the largest
gains at elevated RH. The second result is the more interesting one: \textbf{restricting
operation to 20\,min/day measurably reduced long-term drift} relative to continuous
running, implicating cumulative optical-chamber exposure rather than elapsed calendar
time as the ageing mechanism. Single-site, single sensor model, single season, so the
polynomial coefficients are not transferable --- but duty-cycling is a design lever, not a
calibration constant, and it applies to any continuously powered node.}
{10.1080/26395940.2026.2616098}

\paperentry{Qian et al.\ 2026}
{Enhancing accuracy of indoor air-quality sensors via automated machine-learning
calibration}
{Atmos Meas Tech}
{Multi-stage AutoML calibration chaining field sensors to intermediate drift-correction
reference sensors to a reference-grade instrument, with \textbf{separate models fitted
below and above a concentration break} for clean air versus pollution events. Evaluated
in a controlled indoor chamber with two sensor models against diverse real indoor sources
under uncontrolled ambient conditions: \textbf{$R^2 > 0.90$, NRMSE and sMAPE roughly
halved} against uncalibrated output, bias effectively removed. The split-range design is
the transferable idea --- a single global calibration is being asked to serve two distinct
optical regimes and it fails at both ends. Chamber conditions make the result an upper
bound on field performance, and no ageing term is included, which the Kumpika result
suggests is the dominant long-run error.}
{10.5194/amt-19-603-2026}

\paperentry{Le et al.\ 2026}
{Operational performance of the Vaisala CL61 ceilometer for atmospheric profiling}
{Atmos Meas Tech}
{Three years of operational data from multiple CL61 units --- a depolarisation-capable
elastic backscatter lidar --- assessed for inter-unit consistency, drift and bias. Most
early production units showed \textbf{progressive laser-power decline with rising
background noise}, and manual calibration against atmospheric targets at Lindenberg
showed that \textbf{once laser power falls below 40\% the internal calibration factor
stops compensating}. Instrumental background and bias vary with temperature; a correction
and uncertainty-estimation method is given. Under low aerosol load the bias correction
shifts volume linear depolarisation ratio by up to 0.005, while the volume-to-particle
depolarisation difference reaches 0.1 --- so molecular correction, not instrumental bias,
dominates quantitative aerosol interpretation at 910.55\,nm. One unit's signal loss was
traced to window fogging from insufficient firmware-driven heating. Instructive as a
template: silent, firmware-mediated degradation is exactly the failure mode a distributed
network cannot detect without co-location.}
{10.5194/amt-19-4923-2026}

\paperentry{Silva et al.\ 2026}
{Low-cost sensors for indoor air-quality monitoring: systematic review of accuracy,
applications and limitations}
{J Air Waste Manag Assoc}
{PRISMA review, \textbf{24 studies retained from 2{,}736 screened}, covering optical
particle counters and NDIR CO$_2$ against reference-grade instruments in real settings.
Strong correlations are achievable under controlled or semi-controlled conditions, but
performance degrades with humidity, temperature and source type, and the authors flag
methodological heterogeneity, near-absent inter-device validation, regional bias and
\textbf{unaddressed sensor ageing and cross-sensitivity} as systematic gaps. The
operational conclusion --- that most devices do not meet regulatory compliance
requirements and standardised validation protocols do not exist --- is well supported and
worth citing directly. As a review it inherits the heterogeneity it documents, and a
24-study base is thin for the breadth of the claims.}
{10.1080/10962247.2026.2624795}

\paperentry{Nurseitov et al.\ 2026}
{Design, deployment and field evaluation of a low-cost IoT monitoring network for urban
particulate matter}
{Information}
{Six custom nodes (Winsen ZPHS01B on Raspberry Pi Zero 2 W, $\sim$US\$100/node) deployed
across Almaty over a February--April campaign, \textbf{$>$70{,}000 measurements}, with
three engineering contributions: an active airflow-stabilisation enclosure decoupling
sampling from wind, context-aware adaptive edge filtering cutting transmitted volume
$\sim$40\%, and a segmented telemetry pipeline. Cross-validation was against a consumer
monitor with documented traceability rather than a reference analyser
($R^2 = 0.89$--0.95), and the authors state explicitly that the network provides
\emph{internally consistent} rather than reference-equivalent concentrations --- unusually
honest framing. \textbf{Daily mean PM$_{2.5}$ exceeded the WHO 24\,h guideline on 84\% of
days} (February 54.4 vs March 21.9\,$\mu$g/m$^3$); a PM$_{2.5}$/PM$_{10}$ ratio near 0.96
plus higher weekend levels points to residential coal heating. That ratio is high enough
to warrant suspicion of the coarse channel rather than confidence in the source
inference.}
{10.3390/info17070642}

\band{Sensing --- lower-tier records\hfill 2 entries}{Teal}

\paperentry{Reis 2026}
{Low-power embedded sensor node for environmental monitoring with on-board
machine-learning inference}
{Sensors}
{Architecture and system-level optimisation of a multi-parameter node coupling a
low-power microcontroller to a neural inference accelerator, with adaptive duty-cycling,
hierarchical power domains and quantised fixed-point models doing three-class
normal/anomalous/critical classification at the edge. The reported figures ---
\textbf{94\% inference accuracy, 0.87\,ms latency, $\approx$2.9\,mWh average power,
$\approx$88\% reduction in LoRaWAN transmissions} --- come from MATLAB/Simulink
simulation, not deployment, and there is no co-location against a reference instrument
anywhere in the study. Treat the energy budget as a design target and the accuracy figure
as unvalidated. The duty-cycling and hierarchical-power argument is the reusable part,
and it is also the part that does not depend on the classifier working.}
{10.3390/s26020703}

\paperentry{Mohapatra et al.\ 2026}
{IoT--machine-learning proof-of-concept framework for real-time urban particulate matter
prediction}
{IEEE Access}
{PMS7003 on an Arduino module at campus scale, with outlier treatment, normalisation and
temporal indexing feeding random-forest and polynomial regression, plus $k$-means
clustering to identify persistent high-exposure zones. Predictive performance
\textbf{exceeds 85\% for PM$_{1.0}$ and PM$_{2.5}$ but degrades sharply at 5 and
10\,$\mu$m}, which the authors correctly attribute to discretised, low-variability sensor
output rather than to the model --- a useful negative result about the coarse channels of
optical counters generally. No reference co-location, single site, and the models predict
the sensor's own future readings rather than true concentration, so ``performance'' here
is self-consistency. Included for the size-channel finding only.}
{10.1109/access.2026.3651996}

\band{Exposure assessment \& modelling\hfill 3 full entries}{Sage}

\paperentry{Saeedi et al.\ 2026}
{Feature selection and cross-validation strategy in land-use regression for
mobile-monitored ultrafine particles}
{Environ Sci Technol}
{Toronto mobile UFP campaign used to compare LUR models trained under random, spatial,
temporal and spatiotemporal cross-validation, with and without forward feature selection,
each evaluated on a hold-out test set and against \emph{independent} stationary backyard
measurements. \textbf{Average percentage error fell from $\sim$217\% under random CV to
$\sim$79\% under spatiotemporal CV with FFS.} Random-CV models overfitted, were sensitive
to outliers, produced spatial artefacts and exaggerated variable effects, and failed on
independent samples --- while reporting good internal validation statistics. Spatial folds
were designed at scales reflecting UFP autocorrelation, which is the methodological core.
The blunt implication: a large share of published LUR performance numbers are measuring
data reproduction, and any exposure surface fed into an epidemiological model should
carry a spatiotemporally validated error, not a random-CV one.}
{10.1021/acs.est.5c12601}

\paperentry{Blanco et al.\ 2026}
{Influence of on-road mobile monitoring design on ultrafine particle exposure models and
cognitive health inferences}
{J Expo Sci Environ Epidemiol}
{A reference roadside UFP model built from repeated measurements at 309 locations was
used with the Adult Changes in Thought cohort ($n = 5{,}283$) to test late-life cognitive
function, then \textbf{600\,km of on-road measurements were subsampled} to emulate
alternative campaign designs varying visit frequency, spatial balancing and sampling
time, with universal kriging plus partial least squares. The reference model gave a null:
\textbf{CASI-IRT $+0.007$ (95\% CI $-0.013$ to $0.027$) per 1{,}900\,pt/cm$^3$}.
Alternative designs reproduced the null but attenuated it; route-based sampling and very
short 4-visit campaigns produced the most variable estimates, while temporal and plume
adjustments changed little. Testing design sensitivity against a null limits what can be
concluded --- attenuation toward zero is unobservable when the reference is zero --- and
the authors say so. The design guidance, that visit count and out-of-hours coverage
matter more than post-hoc adjustment, is what carries.}
{10.1038/s41370-026-00845-y}

\paperentry{Zhou M et al.\ 2026}
{Surface PM$_{2.5}$ in 2022 India: emission updates, WRF-Chem evaluation and source
attribution}
{Atmos Chem Phys}
{First WRF-Chem evaluation and source attribution for India for 2022, with an updated
residential inventory reflecting the household clean-fuel transition, a plant-level
coal-power inventory, a revised secondary-organic-aerosol scheme and improved
near-surface mixing. \textbf{Annual PM$_{2.5}$ bias $0.2 \pm 16.9$\,$\mu$g/m$^3$
($0 \pm 31\%$) across 288 South Asian monitoring sites}; simulated population-weighted
annual mean \textbf{47.4\,$\mu$g/m$^3$}. Attribution: transboundary transport 27\%
(12.8\,$\mu$g/m$^3$) --- larger than any single domestic sector --- industry 18\%
(8.6\,$\mu$g/m$^3$), residential 15\% (7.3, still dominant in the Indo-Gangetic Plain),
power 13\% (6.1, concentrated in central India). A $\pm 31\%$ annual bias is respectable
for a CTM but is the same order as the sectoral shares being separated, so the ranking of
industry above residential should be read as provisional. The residential demotion is
nonetheless a policy-relevant result about a transition that actually happened.}
{10.5194/acp-26-10533-2026}

\clearpage

\band{Neuro, cognition \& mental health\hfill 2 full entries}{Deep}

\paperentry{Zhang et al.\ 2026}
{PM$_{2.5}$ and its components and cognitive dysfunction among rural older adults}
{Neurotoxicology}
{\textbf{14{,}379 Henan Rural Cohort participants aged $\geq$60}, cognitive dysfunction
defined by MMSE and the Hasegawa Dementia Scale, with PM$_{2.5}$, black carbon, organic
matter, ammonium, sulfate and nitrate from the TAP dataset; logistic and weighted
quantile sum regression, then network toxicology and random forest for candidate genes.
\textbf{Per-IQR increases raised CD odds by 19--43\% across the six exposures}, and
\textbf{BC was the leading mixture contributor under all four CD definitions} (weights
54--86\%), with OM second (7--35\%). Cross-sectional, so reverse causation through
cognitively impaired participants' residential and behavioural patterns is unaddressed,
and component fields from a single reanalysis product are strongly collinear --- WQS
weights under collinearity are unstable. The BC ranking is consistent with the emerging
component literature; the network-toxicology arm is hypothesis-generating and labelled as
such by the authors.}
{10.1016/j.neuro.2026.103539}

\paperentry{Daoud et al.\ 2026}
{Air pollution and suicide mortality in the UK}
{Environ Sci Pollut Res Int}
{Balanced panel of \textbf{$>$2{,}200 local-authority-year observations, 2015--2023},
linking ONS suicide data to DEFRA modelled PM$_{2.5}$ with local-authority and year fixed
effects, plus nonlinear specifications. The raw cross-sectional association is
\emph{negative}; after fixed effects it reverses to \textbf{$\beta = 0.221$
($p = 0.007$)}, with a U-shaped form turning at $\approx$7.8\,$\mu$g/m$^3$. That sign
reversal is the whole result and it is not reassuring: the estimate is driven entirely by
within-area variation over a period of falling concentrations and rising suicide rates,
with no time-varying socioeconomic controls (deprivation, unemployment, service
provision) in the model. Ecological, modelled exposure, no individual-level confounders.
The authors describe it as exploratory, which is the correct reading. Logged here because
a second, independent PM--suicide record surfaced in the same window (Cai et al., below).}
{10.1007/s11356-026-38089-w}

\band{Reproductive \& developmental\hfill 2 full entries}{Sky}

\paperentry{Pai et al.\ 2026}
{Household incense burning and small-for-gestational-age risk}
{Pediatr Neonatol}
{\textbf{18{,}666 mother--infant pairs} in the Taiwan Birth Cohort Study, incense exposure
during pregnancy categorised none / occasional / daily, SGA stratified by term and
preterm. SGA prevalence rose monotonically across exposure groups (8.1\%, 9.4\%, 10.0\%).
Among term infants, \textbf{aOR 1.16 (1.00--1.33) occasional and 1.25 (1.10--1.41)
daily}; preterm infants showed the same direction without significance, and there was no
association with preterm birth itself. Sensitivity analyses covered other indoor
environmental factors, maternal medical history, gestational complications, gestational
weight gain and secondhand-smoke interaction. The limitation that matters is that
exposure is a self-reported frequency category with no concentration, duration or
ventilation information --- the dose--response is ordinal in \emph{practice}, not in dose,
and the contrast between an unventilated daily burn and a well-ventilated one sits
entirely inside the top category. Pairs directly with Zhou L et al.\ in this issue.}
{10.1016/j.pedneo.2026.06.003}

\paperentry{Bello et al.\ 2026}
{Prenatal PM$_{2.5}$ exposure and infant health in Quebec}
{Econ Hum Biol}
{Administrative birth records for Quebec 2008--2015 linked to monitoring data and tax
records, in neighbourhood-by-quarter fixed-effects models --- a low-pollution setting,
which is the point. \textbf{Population-average effects on birth outcomes were essentially
null}; the result lives entirely in the heterogeneity. Among families in the most
vulnerable stratum of a cumulative socioeconomic vulnerability index, a
\textbf{1\,$\mu$g/m$^3$ increase in mean PM$_{2.5}$ was associated with $\approx$6.6\,g
lower birth weight}, with further effects on prematurity and birth weight among newborns
admitted to neonatal care. \textbf{Days exceeding the Canadian 24\,h standard had
disproportionately large effects in vulnerable families}, so peaks matter beyond the
mean --- directly relevant to any monitoring design optimised for annual averages. No
confidence intervals are reported in the abstract, and the vulnerability index is
constructed rather than pre-registered, which invites specification search across strata.}
{10.1016/j.ehb.2026.101631}

\band{Occupational \& indoor exposure\hfill 2 full entries}{Clay}

\paperentry{Zhou L et al.\ 2026}
{Indoor burning of Arabian incense generates abundant ultrafine particles with strong
oxidative potential}
{Atmos Chem Phys}
{Chamber characterisation of Bakhoor under charcoal-assisted heating representative of
household practice, with acellular DTT activity and a macrophage intracellular
oxidative-stress assay, benchmarked against protocol-matched sidestream cigarette smoke.
Normalised per gram burned: \textbf{670--1690\,$\mu$g\,min$^{-1}$\,g$^{-1}$ mass and
(6--7)$\times$10$^{11}$ particles\,min$^{-1}$\,g$^{-1}$ number, with ultrafines
70--75\% of number} and UFP emission rates above sidestream cigarette smoke. Emission
magnitude followed a nonlinear power law in the hexane-soluble fraction, identifying a
compositional control. \textbf{Water-soluble OP$_{\mathrm{DTT}}$ was
$\approx$32\,pmol\,min$^{-1}$\,$\mu$g$^{-1}$ --- modestly below cigarette particles ---
yet the cellular response was stronger}, and ozone ageing raised OP for both sources at
ageing equivalent to days of indoor residence. The acellular/cellular divergence is the
finding with teeth: DTT alone would have ranked this source as the safer of the two.
Chamber conditions and a charcoal co-source mean the absolute rates are not a household
exposure estimate.}
{10.5194/acp-26-10587-2026}

\paperentry{van Rheenen et al.\ 2026}
{Microbial contamination of air-conditioning and humidifying systems in hospitals and the
built environment}
{J Hosp Infect}
{Mapping review of pathogen presence, biofilm formation and dissemination in the air and
water compartments of air-handling systems analogous to adiabatic humidifiers, searched
across three databases: \textbf{22 of 468 studies retained, covering 184 air-handling
units, none of which used adiabatic humidification}. \emph{Pseudomonas aeruginosa} and
\emph{Burkholderia cepacia} appeared occasionally in system air; \emph{Mycobacterium} and
\emph{Acinetobacter} on cooler surfaces and in drain pans; \emph{Mucor} and
\emph{Aspergillus fumigatus} in water and surrounding air, \emph{Fusarium} more airborne.
\textbf{No included study examined biofilm} --- the proposed mechanism is therefore
unevidenced within the review's own base. The conclusion is a statement about absence of
evidence for the technology actually being adopted on energy-efficiency grounds, which is
a legitimate thing to publish. Bioaerosol rather than mass concentration, so it sits at
the edge of this watch's scope; included because humidification is a growing indoor-air
intervention with no monitoring standard.}
{10.1016/j.jhin.2026.07.020}

\band{Mechanistic toxicology\hfill 1 full entry}{Violet}

\paperentry{Mishra et al.\ 2026}
{Aerosol oxidative potential and reactive species predicted with a chemical kinetics
model (KM-OP)}
{Atmos Chem Phys}
{A chemical kinetic model of oxidative potential built on literature rate coefficients and
a large compilation of laboratory data for transition-metal ions, quinones and organic
aerosol, predicting ROS production and consumption of ascorbic acid and DTT from PM
composition. Applied to field measurements from \textbf{Grenoble, Paris and London:
$R^2 > 0.75$ with $<15\%$ model bias in 4 of 6 datasets}. \textbf{Organic aerosol is the
strongest contributor to antioxidant-based OP assays}, with minor Cu and Fe
contributions; Cu dominates H$_2$O$_2$ production but barely affects \textbullet OH,
while the antioxidant assays correlate well with modelled \textbullet OH --- separating
the species that drives the assay from the species that drives the metal signal. Two of
six datasets are not reproduced, and the model is calibrated on the same body of
literature it is tested against in kind. Even so, this converts OP from a per-filter
measurement into a quantity predictable from speciated composition, which is exactly what
a composition-resolving sensor network would need to be worth deploying.}
{10.5194/acp-26-10605-2026}

\vspace{4mm}

% ---------------------------------------------------------- metadata-only tier
\section*{\sffamily\bfseries\color{Deep}Surfaced without abstract --- metadata only}
\vspace{-1.5mm}{\color{Amber}\rule{\linewidth}{1.1pt}}\vspace{2mm}

{\fontsize{7.4}{9}\selectfont\color{Slate}Eight records reached this issue through the
Crossref by-journal leg but have no abstract deposited in Crossref and are not yet in
Europe PMC or PubMed. They are recorded here on title, journal, deposit date and DOI
alone. \textbf{No summary is offered because none can be written without inventing
content}; several are clearly relevant and should be revisited once abstracts appear.\par}

\vspace{2mm}
\footnotesize
\begin{itemize}\setlength{\itemsep}{1.4mm}
\item \textbf{Cai et al.} --- PM$_{2.5}$ as the dominant air pollutant associated with
suicide among people living with HIV/AIDS: a province-wide case-crossover study in China.
\textit{Atmos Environ}, deposited 29 July.
\texttt{10.1016/\allowbreak j.atmosenv.\allowbreak 2026.122271}
\;{\color{Slate}[pairs with Daoud et al.]}

\item \textbf{Lu et al.} --- Antagonistic effects of PM$_{2.5}$ and ozone on respiratory
morbidity: multi-pollutant interaction analysis using outpatient data from the Pearl
River Delta. \textit{Atmos Environ}, deposited 30 July.
\texttt{10.1016/\allowbreak j.atmosenv.\allowbreak 2026.122270}
\;{\color{Slate}[antagonism, not synergy]}

\item \textbf{Zheng et al.} --- Augmented transformer framework for quantification of
PM$_{2.5}$ chemical composition via one-step laser desorption/ionisation single-particle
aerosol mass spectrometry. \textit{Atmos Environ}, deposited 28 July.
\texttt{10.1016/\allowbreak j.atmosenv.\allowbreak 2026.122260}
\;{\color{Slate}[speciation from SPAMS --- directly on-topic]}

\item \textbf{Park et al.} --- Diagnosing explainability of machine learning:
transformer-based PM$_{2.5}$ estimation in Seoul. \textit{Atmos Environ}, deposited
28 July. \texttt{10.1016/\allowbreak j.atmosenv.\allowbreak 2026.122264}

\item \textbf{Delgado Peralta et al.} --- Improving PM$_{2.5}$ simulations in
metropolitan S\~{a}o Paulo through PMF-based emission-inventory adjustment and local road
transport emissions. \textit{Atmos Environ}, deposited 27 July.
\texttt{10.1016/\allowbreak j.atmosenv.\allowbreak 2026.122249}

\item \textbf{Faggi et al.} --- Heatwaves and particulate matter increases: an Italian
case study. \textit{Atmos Environ}, deposited 27 July.
\texttt{10.1016/\allowbreak j.atmosenv.\allowbreak 2026.122246}

\item \textbf{Yu et al.} --- Indoor particulate matter and associated microplastic
exposure in the educational environment of Shenzhen. \textit{Atmos Pollut Res},
deposited 30 July. \texttt{10.1016/\allowbreak j.apr.\allowbreak 2026.103162}

\item \textbf{Saygin et al.} --- Source identification and seasonal changes of trace
metals, protein content and oxidative potential of fine and coarse particulates from
Sanliurfa after the Kahramanmaras earthquakes. \textit{Atmos Pollut Res}, deposited
31 July. \texttt{10.1016/\allowbreak j.apr.\allowbreak 2026.103164}
\end{itemize}

\vspace{4mm}

% ==================================================================== TABLE
\section*{\sffamily\bfseries\color{Deep}Corpus register}
\vspace{-1.5mm}{\color{Amber}\rule{\linewidth}{1.1pt}}\vspace{2.5mm}

{\fontsize{6.7}{8.4}\selectfont
\setlength{\tabcolsep}{2.4pt}
\begin{longtable}{@{}>{\RaggedRight}p{25mm}>{\RaggedRight}p{41mm}>{\RaggedRight}p{27mm}%
>{\RaggedRight}p{23mm}>{\RaggedRight}p{18mm}>{\RaggedRight}p{15mm}@{}}
\toprule
\textsf{\textbf{First author}} & \textsf{\textbf{Focus}} & \textsf{\textbf{Journal}} &
\textsf{\textbf{Design}} & \textsf{\textbf{Metric}} & \textsf{\textbf{Region}}\\
\midrule
\endfirsthead
\toprule
\textsf{\textbf{First author}} & \textsf{\textbf{Focus}} & \textsf{\textbf{Journal}} &
\textsf{\textbf{Design}} & \textsf{\textbf{Metric}} & \textsf{\textbf{Region}}\\
\midrule
\endhead
\multicolumn{6}{@{}l@{}}{\cellcolor{Mist}\textsf{\textbf{\textcolor{Deep}{Neuro / mental health}}}}\\[0.6mm]
\href{https://doi.org/10.1001/jamanetworkopen.2026.25015}{Dales} & Wildfire-sourced PM$_{2.5}$ and ED visits for migraine and other primary headache & \textit{JAMA Netw Open} & Case-crossover & PM$_{2.5}$ (wildfire) & Canada\\
\href{https://doi.org/10.1093/ageing/afag215}{Glover} & Umbrella meta-review of urban environmental exposures and cognitive health & \textit{Age Ageing} & Umbrella review / meta-analysis & PM$_{2.5}$ & Global\\
\href{https://doi.org/10.1016/j.nsa.2026.107019}{Goossens} & In utero air pollution exposure and pre-/postnatal brain development & \textit{Neurosci Appl} & Narrative review & PM$_{2.5}$ / PM$_{10}$ / PM$_{\mathrm{coarse}}$ & Global\\
\href{https://doi.org/10.1016/j.isci.2026.116617}{Ma} & PM$_{2.5}$-induced neurotoxicity and depression: exposure risks and mechanisms & \textit{iScience} & Narrative review & PM$_{2.5}$ + components & Global\\
\href{https://doi.org/10.3760/cma.j.cn112150-20250804-00752}{Pan} & Maternal PM exposure and offspring neurodevelopment: effects and mechanisms & \textit{Zhonghua Yu Fang Yi Xue Za Zhi} & Narrative review & PM (bulk) & China\\
\addlinespace[0.9mm]
\multicolumn{6}{@{}l@{}}{\cellcolor{Mist}\textsf{\textbf{\textcolor{Teal}{Cardiovascular \& metabolic}}}}\\[0.6mm]
\href{https://doi.org/10.3760/cma.j.cn112338-20251120-00843}{Huang} & PM$_{2.5}$ and dyslipidaemia in rural Chinese children; dietary effect modification & \textit{Zhonghua Liu Xing Bing Xue Za Zhi} & Cross-sectional (multistage) & PM$_{2.5}$ & China\\
\href{https://doi.org/10.32394/pe/218966}{Kazmierska} & Air pollution and cardiovascular disease: review of current data & \textit{Przegl Epidemiol} & Narrative review & PM$_{2.5}$ / PM$_{10}$ & Global\\
\href{https://doi.org/10.1007/s00125-026-06814-2}{Xi} & Air pollution, built environment and temperature vs. prediabetes phenotypes & \textit{Diabetologia} & Prospective cohort & PM$_{2.5}$ & Germany\\
\href{https://doi.org/10.1161/JAHA.125.046917}{Yang} & Short-term PM$_{2.5}$ and hospitalisation/ED use in atrial fibrillation & \textit{J Am Heart Assoc} & Case-crossover & PM$_{2.5}$ (wildfire-influenced) & USA\\
\addlinespace[0.9mm]
\multicolumn{6}{@{}l@{}}{\cellcolor{Mist}\textsf{\textbf{\textcolor{Sky}{Reproductive \& developmental}}}}\\[0.6mm]
\href{https://doi.org/10.1016/j.jes.2025.10.043}{Chen} & Prenatal air pollution, fetal growth trajectory phenotypes and child neurodevelopment & \textit{J Environ Sci (China)} & Prospective cohort & PM$_{2.5}$ / PM$_{10}$ & China\\
\href{https://doi.org/10.1016/j.ecoenv.2026.120566}{Huang} & Joint prenatal bisphenol + air-pollutant mixture exposure and preschool IQ & \textit{Ecotoxicol Environ Saf} & Prospective cohort (mixtures) & PM$_{2.5}$ / PM$_{10}$ & China\\
\href{https://doi.org/10.1007/s40572-026-00552-8}{Peterson} & Ambient air pollution, PAHs and pubertal development: critical review & \textit{Curr Environ Health Rep} & Critical review & PM$_{2.5}$ components & Global\\
\href{https://doi.org/10.1016/j.envres.2026.125313}{Xie} & Long-term PM$_{2.5}$ and timing of pubertal development in boys and girls & \textit{Environ Res} & Prospective cohort & PM$_{2.5}$ & China\\
\addlinespace[0.9mm]
\multicolumn{6}{@{}l@{}}{\cellcolor{Mist}\textsf{\textbf{\textcolor{Sage}{Respiratory \& allergic}}}}\\[0.6mm]
\href{https://doi.org/10.1007/s00484-026-03289-z}{Bayar Muluk} & Climate change and otolaryngology: airway disease impacts and emission reduction & \textit{Int J Biometeorol} & Narrative review & PM$_{2.5}$ & Global\\
\href{https://doi.org/10.4103/mgr.MEDGASRES-D-25-00164}{Liao} & Short-term air pollution, respiratory symptoms and lung function in asthma & \textit{Med Gas Res} & Panel study & PM$_{2.5}$ / PM$_{10}$ & China\\
\href{https://doi.org/10.3390/epidemiologia7040097}{Sanchez de Prada} & Ambient PM and asthma-coded admissions in children 0-3 y, Spain & \textit{Epidemiologia} & Ecological panel & PM$_{2.5}$ / PM$_{10}$ & Spain\\
\href{https://doi.org/10.1016/j.jhazmat.2026.143076}{Yang} & Allergenic vs. toxic risk of airborne seed hairs compared with ambient PM$_{2.5}$ & \textit{J Hazard Mater} & Field measurement & PM$_{2.5}$ + bioaerosol & China\\
\addlinespace[0.9mm]
\multicolumn{6}{@{}l@{}}{\cellcolor{Mist}\textsf{\textbf{\textcolor{Violet}{Mechanistic toxicology}}}}\\[0.6mm]
\href{https://doi.org/10.1016/j.jes.2025.12.074}{Ding} & PM$_{2.5}$ promotes vascular calcification via miR-27b-5p/HES1 signalling & \textit{J Environ Sci (China)} & In vitro + organoid + murine & PM$_{2.5}$ & China\\
\href{https://doi.org/10.1007/s00360-026-01698-5}{Maraschi} & Osmoregulatory disruption and metal bioaccumulation from settleable atmospheric PM & \textit{J Comp Physiol B} & Controlled exposure (ecotox) & Settleable APM & Brazil\\
\href{https://doi.org/10.1016/j.jes.2025.11.018}{Wu} & Macrophage NLRP3 inflammasome in real-ambient PM-induced hepatic glucose dysregulation & \textit{J Environ Sci (China)} & Real-ambient murine + co-culture & Real-ambient PM & China\\
\addlinespace[0.9mm]
\multicolumn{6}{@{}l@{}}{\cellcolor{Mist}\textsf{\textbf{\textcolor{Clay}{Occupational \& indoor}}}}\\[0.6mm]
\href{https://doi.org/10.3389/fmedt.2026.1796828}{Chen} & Nested suction instrument to cut intraoperative surgical-smoke PM$_{2.5}$ & \textit{Front Med Technol} & Bench / device experiment & PM$_{2.5}$ (surgical smoke) & China\\
\href{https://doi.org/10.1093/annweh/wxag055}{Cole} & Determinants of respirable crystalline silica exposure in Australian tunnelling & \textit{Ann Work Expo Health} & Observational exposure study & Respirable dust / RCS & Australia\\
\href{https://doi.org/10.1016/j.jhazmat.2026.143081}{Ding} & Stage-resolved hair Zn-isotope dynamics as a firefighting combustion-exposure tracer & \textit{J Hazard Mater} & Repeated-measures biomarker & Combustion-derived PM & China\\
\href{https://doi.org/10.3390/jox16040119}{Ramsden} & Xenobiotic hazards in aircraft cabin air & \textit{J Xenobiot} & Review + risk economics & Ultrafine particles & Global\\
\addlinespace[0.9mm]
\multicolumn{6}{@{}l@{}}{\cellcolor{Mist}\textsf{\textbf{\textcolor{Amber}{Sensing, forecasting \& instrumentation}}}}\\[0.6mm]
\href{https://doi.org/10.1016/j.jenvman.2026.130515}{Shi} & Adaptive spatiotemporal graph transformer for multi-site multi-step PM$_{2.5}$ forecasting & \textit{J Environ Manage} & Deep-learning model & PM$_{2.5}$ & China / India\\
\href{https://doi.org/10.3390/bios16070380}{Tsuchiya} & Seasonal and indoor-outdoor characteristics of fluorescent aerosol particles in offices & \textit{Biosensors} & Real-time sensor campaign & Bioaerosol / coarse & Japan\\
\addlinespace[0.9mm]
\multicolumn{6}{@{}l@{}}{\cellcolor{Mist}\textsf{\textbf{\textcolor{Amber}{Exposure assessment \& modelling}}}}\\[0.6mm]
\href{https://doi.org/10.1029/2025GH001621}{Solorzano-Araque} & Multiscale exposure modelling: LOTOS-EUROS + mobility surveys + particle toxicity & \textit{GeoHealth} & CTM + mobility integration & PM$_{2.5}$ / PM$_{10}$ & Colombia\\
\addlinespace[0.9mm]
\multicolumn{6}{@{}l@{}}{\cellcolor{Mist}\textsf{\textbf{\textcolor{Coral}{Burden, policy \& mitigation}}}}\\[0.6mm]
\href{https://doi.org/10.1002/hsr2.72886}{Karamzad} & Ischaemic heart disease burden in MENA 1990-2021; PM attributable fraction 28.2\% & \textit{Health Sci Rep} & GBD secondary analysis & Ambient PM & MENA\\
\href{https://doi.org/10.1371/journal.pone.0347989}{Rees} & HEPA filtration to reduce respiratory infection in care homes & \textit{PLoS One} & RCT process evaluation & Filtered PM & UK\\
\href{https://doi.org/10.1016/j.jes.2026.01.009}{Shan} & Policy-driven emission reductions and post-pandemic rebound, high-resolution inventories & \textit{J Environ Sci (China)} & Emission inventory & PM emissions & China\\
\addlinespace[0.9mm]
\multicolumn{6}{@{}l@{}}{\cellcolor{Mist}\textsf{\textbf{\textcolor{Slate}{Other clinical endpoints}}}}\\[0.6mm]
\href{https://doi.org/10.3390/dj14070440}{Lupita} & PM$_{2.5}$, preventive-care capacity and caries/edentulism burden across the EU & \textit{Dent J} & Ecological panel & PM$_{2.5}$ & EU-27\\
\href{https://doi.org/10.1016/j.jdin.2026.06.008}{Park} & Long-term PM exposure and risk of androgenetic alopecia & \textit{JAAD Int} & Nationwide cohort & PM$_{2.5}$ / PM$_{10}$ & South Korea\\
\href{https://doi.org/10.3171/2026.3.PEDS2615}{Sun} & Environmental exposure and area deprivation vs. paediatric hydrocephalus outcomes & \textit{J Neurosurg Pediatr} & Registry cohort + geospatial & PM$_{2.5}$ (satellite-derived) & USA\\
\addlinespace[0.9mm]
\bottomrule
\end{longtable}}


\vspace{4mm}
\begin{keybox}
{\sffamily\bfseries\footnotesize\color{Deep}Provenance and method}\\[1.2mm]
{\fontsize{7.2}{8.8}\selectfont
\textbf{Windows.} The PubMed and Europe PMC legs cover \textbf{1 August 2026} alone,
contiguous with the 31 July issue --- no gap. PubMed was queried on \texttt{[Date -
Entry]} along two independent axes (health, instrumentation): 11 records in the window,
5 retained. Europe PMC \texttt{CREATION\_DATE:[2026-08-01 TO 2026-08-01]} returned 2, both
already published in the 31 July issue and dropped as duplicates. arXiv returned one
record inside the recency screen (Dirichlet-mixture censoring methods), rejected for
having no PM application. OpenAlex date filtering remains behind a paid plan and was not
used.
\\[1.0mm]
\textbf{New this run: the Crossref by-journal leg.} The 31 July issue diagnosed the
sensing blackout and prescribed a Crossref feed keyed on \texttt{created} for the
journals PubMed does not index. That leg is now implemented in \texttt{harvest.py}
(\texttt{crossref\_journal}, \texttt{crossref\_sensing}) over eight ISSNs:
\textit{Atmos.\ Meas.\ Tech.}, \textit{Atmos.\ Chem.\ Phys.},
\textit{Aerosol Sci.\ Technol.}, \textit{Atmos.\ Environ.}, \textit{J.\ Aerosol Sci.},
\textit{Atmos.\ Pollut.\ Res.}, \textit{Aerosol Air Qual.\ Res.},
\textit{Environ.\ Sci.: Atmos.} A \emph{global} \texttt{from-created-date} filter is
useless for entry-date windowing --- continuous re-indexing returns works back to 2007 ---
but scoped to a single journal ISSN it returns that journal's genuinely new deposits,
verified today. Its first run back-swept \textbf{26 July--1 August} and returned
\textbf{55 deposits, 12 PM-relevant}, all 12 carried here (4 with abstracts,
8 metadata-only). They are dated to this issue because 1 August is when the leg first
surfaced them --- the same convention applied to the Europe PMC records in the 31 July
issue. \textbf{The issue's window is therefore not uniform across sources}, which is
stated rather than hidden.
\\[1.0mm]
\textbf{Consensus leg.} Eight further sensing and exposure-modelling records came from
Consensus sweeps on calibration and drift, co-location against reference monitors,
network siting, edge inference and mobile monitoring. Consensus has no entry-date filter,
so these are screened against \texttt{state/seen.json} rather than by date; two hits
(Ding, Nguyen) were dropped as already published on 29 July. DOIs were recovered by
Crossref bibliographic matching, since Consensus returns none.
\\[1.0mm]
\textbf{Screening.} 25 records retained, \textbf{15 rejected with a logged reason} in
\texttt{state/rejected.jsonl}: 4 duplicates across sources, 3 gas-phase-only
instrumentation, 2 non-atmospheric particulates (pharmaceutical surface cleanliness,
marine sediment resuspension), 1 catalysis review, 1 heat-mortality study in which PM
enters only as a covariate, 1 therapeutic-nebuliser deposition study, 1 generic
statistical method with no PM application, 1 with no retrievable abstract from any
connector, 1 formaldehyde hematotoxicity review.
\\[1.0mm]
\textbf{Trial watch.} The ClinicalTrials.gov sweep on filtration interventions added
three sham-controlled trials to \texttt{state/trials.json} (NCT05867381, NCT05718245,
NCT05016271); they are surfaced in the weekly rollup, not here.
\\[1.0mm]
Subtopic, design, particle-metric and geography labels were assigned from title,
abstract, keywords and MeSH terms; assignment is single-label by dominant contribution
and is a judgement, not a controlled vocabulary. Effect estimates in the forest plot are
transcribed verbatim from abstracts. Every DOI was verified and resolved before
publication. Content derived from PubMed, Europe PMC and Crossref metadata; abstracts
remain the copyright of their respective publishers.\par}
\end{keybox}

\end{document}
